\section{Introduction}
\label{sec:intro}

Long-running LLM agents keep a memory of past interactions: summaries that are periodically rewritten, extracted fact
lists, or retrieval indexes over raw turns \citep{packer2023memgpt,chhikara2025mem0,rasmussen2025zep}. Memory systems
are compared on benchmarks that feed a long history and then ask questions about it
\citep{maharana2024locomo,wu2025longmemeval,hu2025memoryagentbench}. A correct answer is read as evidence that the memory delivered the relevant fact: wrote it, kept it through later
updates, and retrieved it when asked.

That reading is unsafe. Whichever stage fails, whether the fact was never written, was overwritten, or was stored but
not retrieved, the reader model can still answer from what it already believes. A summary that kept nothing about a long document still lets the model answer that Normandy is in France, and a memory
that kept nothing about a user still lets it guess that someone who volunteers for a support group might study
psychology, which is the gold answer in LoCoMo. In our controlled study a widely used compaction prompt kept \aiderRet{} of the facts it was given, yet
its agent answered \aiderCongAcc{} of the questions about schema-congruent facts correctly (\S\ref{sec:controlled}). Raw accuracy
cannot tell this apart from remembering. The prior also settles conflicts: when a user's value changes away from what the
model expects, an agent whose memory lost the update falls back to the old, expected value, while a change toward the
expectation is answered correctly from the prior alone.

The worry is not new in other areas of NLP: question-only and hypothesis-only baselines exposed benchmarks that could
be solved without the evidence \citep{kaushik2018much,poliak2018hypothesis,gururangan2018annotation}, and recent memory
benchmarks build fictional worlds or filter guessable items when they are created
\citep{groundtruthfirst2026,memgym2026,longmemevalvtwo2026}. What is missing is a way to measure, for the benchmarks
the field already reports, how much of a memory score is the prior; and a scoring procedure that removes it with
explicit, testable assumptions. Because a guess can rescue a failure at any stage, such a procedure has to work from final
answers alone, whatever the memory architecture. This paper provides both.

\paragraph{Contributions.}
\begin{itemize}
\item \textbf{An audit} of four agent-memory benchmarks with two reader models: the expected closed-book accuracy of the
reader ranges from \gFcMhPct{} to \gEventqaPct{} across task families (\S\ref{sec:audit}).
\item \textbf{\pcme{}}, a prior-controlled evaluation procedure (\S\ref{sec:formulation}--\S\ref{sec:procedure}). It
defines a primary memory score that needs only one assumption (accuracy on prior-resistant items, \cpr{}), a
guess-corrected score for the full item set that is identified only under a forced-answer protocol, and a
discrimination score for unanswerable items. We validate every metric with pre-registered property tests on rule-built
reference systems and report the two tests that failed in the first version and how the procedure was changed.
\item \textbf{A re-scoring} of twelve memory systems under three answer protocols (\S\ref{sec:audit}). Abstention, not
memory, explains most raw-accuracy differences among memory-poor systems; standardising the protocol alone does not
recover memory-based rankings; and LoCoMo's unanswerable category rewards forgetting.
\item \textbf{A controlled instrument}: \nFinal{} same-slot counter-schema minimal pairs written for this study and
filtered by pre-registered gates on four model families, and a replication of the mechanism with two compactor
families (\S\ref{sec:controlled}).
\end{itemize}
The audit and the controlled instrument answer complementary questions. The audit shows that scores on the benchmarks in
use are inflated and confounded, but those benchmarks cannot reveal what each memory delivered; the controlled instrument
supplies that ground truth, which we use to validate the correction and to locate where the prior enters.
All analyses were pre-registered before the corresponding runs; deviations are listed in Appendix~\ref{app:deviations}.
Code, items, all prompts (verified byte-for-byte against the requests that produced the results) and model outputs
will be released.
